\section{RESULTS AND ANALYSIS}

\subsection{Overall Performance Comparison}

\begin{table}[t]
\caption{Performance comparison with and without TAGA (500 episodes)}
\label{table:main-results}
\centering
\begin{tabular}{lcccccc}
\hline
\textbf{Method} & \textbf{SR↑} & \textbf{CR↓} & \textbf{TR↓} & \textbf{GCR↓} & \textbf{NT(s)} & \textbf{PL(m)} \\
\hline
AG-RL & 0.78 & 0.20 & 0.02 & 10.54\% & 15.47 & 21.54 \\
AG-RL+TAGA & \textbf{0.82} & \textbf{0.16} & \textbf{0.02} & 8.82\% & 16.79 & 22.68 \\
\hline
DS-RNN & 0.64 & 0.32 & 0.04 & 9.87\% & 17.23 & 22.15 \\
DS-RNN+TAGA & 0.66 & 0.29 & 0.05 & 4.13\% & 18.45 & 23.87 \\
\hline
ORCA & 0.59 & 0.35 & 0.06 & 7.31\% & 26.78 & 20.23 \\
ORCA+TAGA & 0.55 & 0.33 & 0.12 & \textbf{4.18\%} & 27.85 & 21.50 \\
\hline
SF & 0.27 & 0.55 & 0.18 & 6.28\% & 29.56 & 21.36 \\
SF+TAGA & 0.29 & 0.53 & 0.18 & 4.22\% & 30.12 & 22.14 \\
\hline
\end{tabular}
\end{table}

Table \ref{table:main-results} demonstrates TAGA's effectiveness across all methods. The key findings are:

\textbf{Group Intrusion Reduction}: TAGA reduces GCR by 16.3\% (AG-RL), 58.2\% (DS-RNN), 42.8\% (ORCA), and 32.8\% (SF). This improvement occurs while GCR remains independent of terminal states, validating our metric design.

\textbf{Safety Enhancement}: AG-RL+TAGA improves both success rate (78\%→82\%) and collision rate (20\%→16\%), demonstrating that the hierarchical safety controller successfully prevents individual collisions during group avoidance maneuvers.

\textbf{Efficiency Trade-off}: TAGA increases navigation time by 1-2 seconds and path length by 1-2 meters—a reasonable cost for socially compliant navigation.

\subsection{Scenario-Based Analysis}

\begin{table}[t]
\caption{TAGA performance across different scenarios (ORCA baseline)}
\label{table:scenarios}
\centering
\begin{tabular}{lccc|ccc}
\hline
 & \multicolumn{3}{c|}{\textbf{ORCA}} & \multicolumn{3}{c}{\textbf{ORCA+TAGA}} \\
\textbf{Scenario} & SR & CR & GCR & SR & CR & GCR \\
\hline
Dense Groups & 0.45 & 0.42 & 12.3\% & 0.52 & 0.38 & 6.7\% \\
Mixed 50-50 & 0.59 & 0.35 & 7.3\% & 0.61 & 0.32 & 4.2\% \\
Dynamic & 0.51 & 0.39 & 8.9\% & 0.54 & 0.36 & 5.1\% \\
Crossing & 0.48 & 0.41 & 9.2\% & 0.53 & 0.37 & 4.8\% \\
Groups Only & 0.43 & 0.45 & 15.2\% & 0.48 & 0.41 & 7.3\% \\
\hline
\end{tabular}
\end{table}

Table \ref{table:scenarios} validates TAGA's robustness across diverse conditions. The most significant improvements occur in dense (45.5\% GCR reduction) and groups-only scenarios (52.0\% GCR reduction) where group awareness is critical. Even in challenging dynamic scenarios with moving groups, TAGA maintains consistent performance improvements.

\subsection{Ablation Study}

\begin{table}[t]
\caption{Effect of safety controller components (AG-RL, 100 episodes)}
\label{table:ablation}
\centering
\begin{tabular}{lccc}
\hline
\textbf{Configuration} & \textbf{SR} & \textbf{CR} & \textbf{GCR} \\
\hline
Baseline AG-RL & 0.78 & 0.20 & 10.54\% \\
TAGA without safety & 0.71 & 0.27 & 7.23\% \\
TAGA with critical zone only & 0.76 & 0.22 & 8.15\% \\
TAGA with full safety (ours) & 0.82 & 0.16 & 8.82\% \\
\hline
\end{tabular}
\end{table}

Table \ref{table:ablation} demonstrates the importance of the hierarchical safety controller. Without it, TAGA reduces group intrusions but increases individual collisions. The full three-zone safety system achieves optimal balance between group avoidance and individual safety.

\subsection{Computational Overhead}
TAGA adds minimal computational cost: group detection (DBSCAN) requires 3.2ms per frame, tangent computation takes 0.8ms per detected group, and safety checking adds 1.5ms. Total overhead of ~5.5ms maintains real-time operation at 10Hz, confirming practical viability.

\subsection{Comparison with Geometric Baselines}
Zone-based and F-formation methods achieved 0\% success rate in our tests due to excessive repulsive forces preventing goal reaching. This fundamental limitation—inability to balance avoidance with goal-directed navigation—highlights TAGA's advantage of computing efficient tangent paths rather than relying on repulsion.

% \vspace{-10pt} 
% \section{RESULTS}
% \vspace{-5pt} 


% \begin{table}[t]
% \caption{Comparison of models with and without TAGA.}
% \label{table:reactive}
% % \caption{An Example of a Table}
% % \label{table_example}
% \begin{center}
% \begin{tabular}{p{2.5cm}cccccc}
%        \hline
%         \textbf{Models} & \textbf{SR$\uparrow$} & \textbf{CR$\downarrow$} & \textbf{GCR$\downarrow$} & \textbf{TR$\downarrow$} & \textbf{NT$\downarrow$} & \textbf{PL$\downarrow$} \\
%         \hline
%         AG-RL\cite{liu2023intention} & \textbf{0.57} & \textbf{0.06} & 0.35 & \textbf{0.02} & \textbf{14.23} & 19.06 \\
%         DS-RNN\cite{liu2020decentralized} & 0.46 & 0.19 & 0.33 & 0.02 & 16.41 & 18.15 \\
%         ORCA\cite{orca} & 0.32 & 0.13 & 0.36 & 0.19 & 24.86 & \textbf{17.65} \\
%         SF\cite{helbing1995social} & 0.26 & 0.39 & 0.14 & 0.21 & 26.79 & 19.49 \\
%         \hline
%         % \textbf{Models + TAGA} & - & - & - & - & - & - \\
%         % \hline
%         AG-RL\cite{liu2023intention}+TAGA & \textbf{0.57} & 0.22 & 0.19 & \textbf{0.02} & 14.61 & 18.83 \\
%         DS-RNN\cite{liu2020decentralized}+TAGA & 0.48 & 0.40 & 0.08 & 0.04 & 15.49 & 18.68 \\
%         ORCA\cite{orca}+TAGA & 0.47 & 0.23 & 0.10 & 0.20 & 27.12 & 19.34 \\
%         SF\cite{helbing1995social}+TAGA & 0.29 & 0.48 & \textbf{0.03} & 0.20 & 27.12 & 19.96 \\
%         \hline 
%   \end{tabular}
% \end{center}
% \vspace{-15pt} 
% \end{table}

% \subsection{Baseline Results}
% \begin{table}[]
%     \centering
%     \caption{Comparison with Existing Models}
%     \label{table:reactive}
%     % \renewcommand{\arraystretch}{1.2} % Adjust row spacing
%   \begin{tabular}{|c|c|c|c|c|c|c}
%        \hline
%         \textbf{Models} & \textbf{Success Rate} & \textbf{Collision Rate} & \textbf{Group Collision} & \textbf{Timeout Rate} & \textbf{Navigation Time (sec)} & \textbf{Path Length (m)} \\
%         \hline
%         Attention-Graph RL~\cite{liu2023intention} & 0.57 & 0.06 & 0.35 & \textbf{0.02} & \textbf{14.23} & 19.06 \\
        
%         DS-RNN~\cite{liu2020decentralized} & 0.54 & 0.24 & 0.18 & 0.04 & 15.26 & 20.39 \\
%         \hline
%         ORCA~\cite{orca} & 0.32 & 0.13 & 0.36 & 0.19 & 24.86 & \textbf{17.65} \\
%         \hline
%         Social Force~\cite{helbing1995social} & 0.26 & 0.38 & 0.14 & 0.16 & 26.79 & 19.49 \\
%         \hline
%         Attention-Graph RL~\cite{liu2023intention} + TAGA & 0.57 & 0.22 & 0.19 & \textbf{0.02} & 14.61 & 18.83 \\
%         \hline
%         DS-RNN~\cite{liu2020decentralized} + TAGA & 0.54 & 0.25 & 0.17 & 0.04 & 15.18 & 20.56 \\
%         \hline
%         ORCA~\cite{orca} + TAGA & 0.47 & 0.23 & 0.10 & 0.20 & 24.86 & 26.79 \\
%         \hline
%         Social Force~\cite{helbing1995social} + TAGA & 0.26 & 0.48 & 0.02 & 0.24 & 26.79 & 19.49 \\
%         \hline 
%   \end{tabular}

%     % \begin{tabular}{|>{\centering\arraybackslash}p{3cm}|>{\centering\arraybackslash}p{2cm}|>{\centering\arraybackslash}p{2cm}|>{\centering\arraybackslash}p{2cm}|>{\centering\arraybackslash}p{2cm}|>{\centering\arraybackslash}p{2cm}|>{\centering\arraybackslash}p{2cm}|}
%     %     \hline
%     %     \textbf{Models} & \textbf{Success Rate} & \textbf{Collision Rate} & \textbf{Group Collision} & \textbf{Timeout Rate} & \textbf{Navigation Time (sec)} & \textbf{Path Length (m)} \\
%     %     \hline
%     %     Attention-Graph RL~\cite{liu2023intention} & 0.57 & 0.06 & 0.35 & \textbf{0.02} & \textbf{14.23} & 19.06 \\
        
%     %     DS-RNN~\cite{liu2020decentralized} & 0.54 & 0.24 & 0.18 & 0.04 & 15.26 & 20.39 \\
%     %     \hline
%     %     ORCA~\cite{den2011a} & 0.32 & 0.13 & 0.36 & 0.19 & 24.86 & \textbf{17.65} \\
%     %     \hline
%     %     Social Force~\cite{helbing1995a} & 0.26 & 0.38 & 0.14 & 0.16 & 26.79 & 19.49 \\
%     %     \hline
%     %     Attention-Graph RL~\cite{liu2023intention} + TAGA & 0.57 & 0.22 & 0.19 & \textbf{0.02} & 14.61 & 18.83 \\
%     %     \hline
%     %     DS-RNN~\cite{liu2020decentralized} + TAGA & 0.54 & 0.25 & 0.17 & 0.04 & 15.18 & 20.56 \\
%     %     \hline
%     %     ORCA~\cite{den2011a} + TAGA & 0.47 & 0.23 & 0.10 & 0.20 & 24.86 & 26.79 \\
%     %     \hline
%     %     Social Force~\cite{helbing1995a} + TAGA & 0.26 & 0.48 & 0.02 & 0.24 & 26.79 & 19.49 \\
%     %     \hline
%     % \end{tabular}
% \end{table}


% \vspace{-5pt} 

% The results in Table \ref{table:reactive} show that TAGA significantly reduces GCR, minimizing collisions with human groups. The reduction for AG-RL is \((0.35 - 0.19)/0.35 \times 100 = 45.7\%\), while DS-RNN, ORCA, and SF achieve reductions of 75.7\%, 72.2\%, and 78.6\%, respectively. This improvement results from incorporating group-reactive attention, which enhances the robot's understanding of group behaviors. Although individual CR increases a bit, it does not affect overall SR. The CR increases because when the TAGA activates, it primarily focuses on group behavior. As a result, if an individual human suddenly appears in the scene at that moment, the robot may fail to react appropriately, leading to a collision. This occurs because the attention mechanism is biased toward group dynamics, reducing responsiveness to unexpected individual obstacles.

% % The results shown in Table \ref{table:reactive} demonstrate that TAGA outperforms existing models in GCR metrics, significantly reducing collisions with human groups. This is achieved through the incorporation of group-reactive attention, which improves the robot's understanding of group behaviors. Although this approach leads to an increase in individual CR but it does not negatively impact the overall SR of the navigation tasks.

% Moreover, the TAGA method exhibits notable improvements in the GCR, highlighting its effectiveness in avoiding groups. However, the CR and TR remain somewhat higher. Additionally, NT and PL are slightly greater compared to the vanilla method. This increase occurs because, while avoiding groups, the robot occasionally chooses longer paths, resulting in a minor rise in both NT and PL.



% \subsection{Simulation Results}
% The following figures in Fig. \ref{fig:sim-res-1} depict key moments during the simulation of the group-aware navigation system. These images show the robot navigating through different crowd scenarios, detecting groups, and avoiding obstacles.
% In the simulation results shown in Fig. \ref{fig:sim-com}, when TAGA detected human groups, it applied tangent actions to keep a safe distance from the group boundaries. This allowed the robot to navigate around the groups without collisions while staying on track towards its goal. As the robot approached the group boundaries, it adjusted its movement to avoid crossing them. Once the robot passed the group boundaries, the tangent action stopped, and the robot switched to its default navigation to reach the target without further interference from the group.

% In the simulation results shown in Fig. \ref{fig:sim-com}, when TAGA detected human groups, it applied tangent actions to maintain a safe distance from the group boundaries, ensuring smooth and collision-free navigation. This mechanism allowed the robot to navigate around groups while staying on track toward its goal. As the robot neared the group boundaries, it dynamically adjusted its trajectory to avoid crossing into the group’s space, prioritizing social compliance. Once the robot successfully passed the group, the tangent action deactivated, allowing it to revert to its default navigation strategy and proceed toward the target without further interference.
% \vspace{-5pt} 
% \subsection{Comparison Results}
% Fig. \ref{fig:sim-com} shows the navigation paths of TAGA (bottom) and AG-RL (top) in the same scenario, highlighting the differences in how each model handles group avoidance and navigation efficiency.

% As shown in Fig. \ref{fig:sim-com}, the AG-RL model does not exhibit group-aware behavior and instead moves directly through the middle of the group. This occurs because AG-RL is primarily designed to find the shortest path to the goal without explicitly considering group formations. It treats all humans as independent obstacles rather than recognizing them as cohesive social entities. Consequently, when multiple humans are clustered together, AG-RL may navigate through them if it perceives an open space. In contrast, TAGA adjusts its trajectory to navigate around groups while maintaining a safe distance, leveraging group-awareness for more socially acceptable and human-like navigation.



% As shown in Fig. \ref{fig:sim-com}, the AG-RL model does not avoid group members, instead passing through the middle of the group. This happens because it tries to find the shortest path to the goal and lacks the ability to avoid groups. In contrast, after detecting the group, TAGA adjusts its path to pass around the group, maintaining a safe distance.